\subsubsection{The Complementing Condition for Elliptic Systems}\label{subsec:elliptic-system-complementing-condition}

For ODE boundary-value problems, it is well known that a ``reasonable\textquotesingle\textquotesingle\ boundary-value problem is obtained if the number of boundary conditions at each point equals half the order of the differential equation and the boundary conditions at each point are linearly independent. For elliptic PDEs, the picture is more complicated. Consider the problem
\begin{equation}\label{eq:complementing-motivation-biharmonic-equation}
\Delta\Delta u=0
\end{equation}
with boundary conditions
\begin{equation}\label{eq:complementing-motivation-biharmonic-boundary}
\Delta u=\frac{\partial}{\partial n}\Delta u=0.
\end{equation}
Although the two boundary conditions are independent of each other, we see that every harmonic function satisfies both the differential equation and the boundary conditions. There are infinitely many linearly independent harmonic functions even within the class of polynomials. This illustrates the need to formulate hypotheses about boundary conditions leading to ``good\textquotesingle\textquotesingle\ problems. These hypotheses will not just express independence of the boundary conditions from each other, but also involve a relationship between the boundary conditions and the differential equation.

The complementing condition provides such a characterization of ``good\textquotesingle\textquotesingle\ boundary conditions. We shall state it for general elliptic systems. As in Chapter~\ref{chap:characteristics}, let us consider a $k\times k$ system of equations
\begin{equation}\label{eq:weighted-elliptic-general-system}
L_{ij}(\bm{x},D)u_j=f_i(\bm{x}),\quad\bm{x}\in\Omega,\ i=1,\ldots,k.
\end{equation}
As before, we assign the ``weights\textquotesingle\textquotesingle\ $s_i$ to the $i$th equation and $t_j$ to the $j$th independent variable in such a way that $L_{ij}$ is at most of order $s_i+t_j$, and we denote the terms that are exactly of order $s_i+t_j$ by $L_{ij}^p$. The condition of ellipticity is that
\begin{equation}\label{eq:weighted-elliptic-system-principal-determinant}
\det\bm{L}^p(\bm{x},\bm{\xi})\ne0\quad\forall\bm{\xi}\in\mathbb{R}^n\backslash\{0\}.
\end{equation}
As we remarked in Chapter~\ref{chap:characteristics}, this condition can be interpreted as follows: Take the values of the coefficients at any fixed point $\bm{x}_0$ in $\Omega$ and consider the system $\bm{L}^p(\bm{x}_0,D)\bm{u}=0$ on all of $\mathbb{R}^n$. Ellipticity means that this constant coefficient system has no nonconstant periodic solutions. The complementing condition will be an analogue of this for points at the boundary.

Let the order of the system, i.e., the order of the polynomial $\det\bm{L}^p(\bm{x},\bm{\xi})$, be $2m$. We impose $m$ boundary conditions
\begin{equation}\label{eq:weighted-elliptic-general-boundary-conditions}
B_{lj}(\bm{x},D)u_j=g_l(\bm{x}),\quad\bm{x}\in\partial\Omega,\ l=1,\ldots,m.
\end{equation}
We now define weights $r_l$ to each boundary condition $l=1,\ldots,m$ so that the order of $B_{lj}$ be bounded by $r_l+t_j$. Again the terms which are precisely of order $r_l+t_j$ will be considered the principal part. (If $r_l+t_j$ is negative, it is of course understood that $B_{lj}=0$.) Let now $\bm{x}_0$ be a point on $\partial\Omega$ and let $\bm{n}$ be the outer normal to $\Omega$. We consider the constant coefficient problem
\begin{equation}\label{eq:complementing-frozen-principal-system}
L_{ij}^p(\bm{x}_0,D)u_j=0,\quad i=1,\ldots,k,
\end{equation}
on the half-space $(\bm{x}-\bm{x}_0)\cdot\bm{n}<0$, with boundary conditions
\begin{equation}\label{eq:complementing-frozen-principal-boundary}
B_{lj}^p(\bm{x}_0,D)u_j=0,\quad l=1,\ldots,m,
\end{equation}
on the boundary $(\bm{x}-\bm{x}_0)\cdot\bm{n}=0$.

\begin{Definition}
\label{def:elliptic-system-complementing-condition}
We say that the \textbf{complementing conditions} holds at $\bm{x}_0$, if there are no nontrivial solutions of \eqref{eq:complementing-frozen-principal-system}, \eqref{eq:complementing-frozen-principal-boundary} of the following form:
\begin{equation}\label{eq:complementing-decaying-tangential-mode}
\bm{u}(\bm{x})=\exp(i\bm{\xi}\cdot(\bm{x}-\bm{x}_0))\bm{v}(\eta),
\end{equation}
where $\bm{\xi}$ is a nonzero real vector perpendicular to $\bm{n}$, $\eta=(\bm{x}-\bm{x}_0)\cdot\bm{n}$, and $\bm{v}(\eta)$ tends to $0$ exponentially as $\eta\to-\infty$.
\end{Definition}
